% =====================================================================
%  Chapter 3 — Principle: Honesty (Hide Nothing)
%  Map: research/chapter-map.md (Chapter 3)
%  Shape: idea → origin → cost → what it opens up. Ideas, not how-to.
%  Rule: the author is never quoted in body text; short quotes only in
%  footnotes. Full statements live in interviews/.
% =====================================================================
\chaptersubtitle{Hide Nothing}
\chapter{Honesty}
\label{ch:hidden}

Most languages describe their goals in terms of speed, safety or
productivity. Silica's public statement of purpose uses a word that is rarely
applied to software at all. It aims to be the most intellectually
\emph{honest} language and runtime in the world, and to get there without
making programmers fight it.\footnote{\src{silica/docs/index.md}: ``the world's
most intellectually honest, secure language and runtime.''} The other nine
principles in this book can each be read as a special case of this one.

\section{What dishonesty looks like}

For Barney, the first dishonesty was not a bug but a reputation. It came with
understanding what compiler optimizers actually do. They remove work that does
nothing and merge work done twice, but most of their effort goes into something
else: translating the behaviour of a PDP-11, the machine C was designed around,
into the behaviour of the very different chips programs really run on. It was
disturbing to realise that, and more disturbing to hear otherwise
knowledgeable people describe the languages that depend on such optimizers as
close to the metal. They were not, and they are not.\footnote{Interview,
2026-11-13.}

For years this was a nagging concern without words. The words arrived in 2018,
when David Chisnall published an article in \emph{ACM Queue} whose title said it
plainly: ``C Is Not a Low-level Language.'' Its subtitle said the rest: ``Your
computer is not a fast PDP-11.''\footnote{D.~Chisnall, ``C Is Not a Low-level
Language: Your computer is not a fast PDP-11,'' \emph{ACM Queue} 16(2), 2018.}
Those ideas stayed at the front of Barney's mind from the day the article
appeared. What was missing, for seven years, were the tools to do anything about
them on a scale that mattered.

% TODO-VERIFY: Chisnall citation details (volume, issue, year).

A language is dishonest, in this sense, when what a program does is not what
the program says. A function that looks like a calculation but writes to a
file. A call that looks local but quietly reserves memory that something else
must clean up later. A library that looks like ordinary code but runs foreign
instructions that can corrupt the whole process. A failure that is caught,
recorded nowhere, and replaced with a plausible answer.

None of these are bugs in the usual sense. They are conventions, and every
mainstream language relies on them. The cost of a convention is that it lives
in people's heads rather than in the program, held together, as Silica's
documents put it, by code reviews and runtime luck.\footnote{\src{silica/docs/index.md}, lines 24--26.} Silica's alternative is to make the program's
own text tell the truth about what it does, and to have the compiler hold it to
that.

% TODO-INTERVIEW: Which language or experience made you name honesty as the
%   target? What was the most dishonest thing a tool ever did to you?

\section{Nothing provided out of sight}

The idea reaches into places most languages treat as plumbing. A Silica
program never receives resources it did not ask for and cannot see: whenever
it starts a new actor, it says how that actor's memory will be managed. Where a
choice could be made either in the program or somewhere outside it, such as a
setting in the environment, the program wins, because a reader of the program
can see it.

\section{No decision made silently}

The same reasoning produces a refusal that surprises newcomers. In several
places Silica has no defaults. A default is a decision someone else made, on
the program's behalf, without saying so. That is convenient in a ten-line
example and corrosive in a large system, where nobody remembers what the
defaults were or why they were chosen.

% TODO-VERIFY: The "no defaults" rule is agent-recorded in .claude-memory;
%   confirm its wording and scope with the author.
% TODO-INTERVIEW: Is there any default you have been tempted to allow?

\section{Honest about limits}

Honesty also applies to what the language cannot do. On a bare-metal board,
Silica can promise things about memory that an operating system will not let it
promise on a laptop, and its documents say so plainly rather than blurring the
difference. Its bridge to foreign code is described as a boundary, explicitly
not a sandbox. Its plainest network connection is documented as offering no
encryption and no identity. Nearly every design document in the project opens
by separating what is built from what is only planned.

Software projects usually describe their hopes in the present tense. Silica's
documents mostly do not, and where they slip, the slip is treated as a defect.

\section{Honest about the tools}

The tools that finally made Silica possible were language models. Barney does
not use them to improvise code on a whim. A colleague's description fits
better: they are an army of developers on Barney's team, working under
direction.\footnote{Interview, 2026-11-13.} The difference in scale is hard to
overstate. Years earlier, getting SwErl to communicate with Erlang machines had
taken well over a month of slow progress. With language models, the Erlang
documentation and Barney's guidance, the equivalent work took about twenty
minutes.

Trust in that army was earned gradually. At first there was no telling how far
machine-written code could be relied on; the hype, for and against, was hard to
separate from the truth without using the models for something significant. The
models and their tooling improved while Silica was being built, and as more of
the compiler existed in Silica, the models had more of it to learn from.

The project's own AI policy has not kept up. It was written to keep AI out of
the authorship of contributions, while much of Silica's history was built with
exactly these tools. Barney regards the policy as out of date. Trust, as Barney
now sees it, belongs with the contributor who uses the tool and answers for
what it produces, not with the tool. How best to say that in a policy is still
being worked out, and this book will not pretend otherwise.

% TODO-INTERVIEW: Revisit when the new AI policy is written.

\section{The rule applied to itself}

In October 2026 the principle caught the project itself. The runtime had been
managing long-lived memory in a way no one had ever decided on, while the
public pages described something else. Barney's reaction was blunt dismay.
Pushing hard for the behaviour the project needed, Barney had not been watching
closely enough what the army had been producing: workers that were highly
focused and, in Barney's word, unscrupulous.\footnote{``my army of highly
focused, yet unscrupulous, workers.'' Interview, 2026-11-13.} It was time to go
back in deeply and fix it.

The response was not to change the documentation to match the runtime. It was
to change the runtime to match the promise (Chapter~\ref{ch:memory}), and to
change how the work is checked. Barney now reviews far more, automates the
reviews where possible, and has the work checked repeatedly against the design
documents.

\section{What it costs}

Honesty is verbose. Programs say more, because they are not allowed to leave
things unsaid. Newcomers meet decisions they would rather not make yet. And a
project that publishes its limits looks, at first glance, less finished than one
that does not.

\section{What it opens up}

A program that tells the truth about itself can be trusted by people who did
not write it. That is a larger opportunity than it first appears. The principle
behind it is simple: the code you see is the code that runs, and code like that
is easier to reason about, to analyse and to debug.

\begin{itemize}
  \item \textbf{Machines can reason about it.} An honest program is easier
        for language models to understand, to search for bugs and to extend.
        Barney now works on the compiler and its sibling projects without an
        integrated development environment or even a conventional editor,
        looking at code directly only when something needs a detailed look.
        There are no tool commands to learn, and the mental load is lighter
        for it.
  \item \textbf{Trouble has an address.} When an application misbehaves,
        the clearly marked trees of \code{dangerous_} code are where the
        problem most probably lies. Reviewers and auditors know where to look
        first.
  \item \textbf{Types in plain sight.} Knowing exactly what kind of data a
        piece of code handles, without unwinding names, is a large advantage
        for anyone who has to check it.
  \item \textbf{Review becomes reading.} A reviewer, an auditor and an
        automated tool all read the same contract in the source, instead of
        reconstructing behaviour from names, comments and experience.
  \item \textbf{Documentation cannot drift.} When behaviour must be stated in
        the program, the most important documentation is checked by the
        compiler every time the program is built.
  \item \textbf{Trust can be delegated.} Code written by a colleague, a
        student or a language model can be judged by what it declares, rather
        than by faith in its author.
  \item \textbf{Hard environments become reachable.} Systems that must be
        certified, audited or defended---medical devices, infrastructure,
        anything regulators care about---need exactly the kind of evidence an
        honest program provides.
\end{itemize}


\begin{sidebox}[Where honesty shows up]
\begin{itemize}\itemsep0pt
  \item Effects stated where they happen (Chapter~\ref{ch:effects}).
  \item Memory that is never provided out of sight (Chapter~\ref{ch:memory}).
  \item Foreign code that cannot hide (Chapter~\ref{ch:danger}).
  \item Failures that are always reported (Chapter~\ref{ch:failsoft}).
  \item Plans never described as finished work.
\end{itemize}
\end{sidebox}
